% preambulo-grego-cardo.tex
% Faz a Cardo entrar sozinha nos trechos em grego, mantendo o resto do
% texto na fonte principal.
% Instalar em: <userdir>\preambulo-grego-cardo.tex
%
% POR QUE ISTO EXISTE
% As fontes desenhadas para acessibilidade — Atkinson Hyperlegible,
% OpenDyslexic — não cobrem grego politônico. Usá-las sem tratamento faria
% o grego SUMIR do PDF, com apenas um "[WARNING] Missing character" na saída.
%
% Este preâmbulo resolve o conflito: o português sai na fonte acessível e
% o grego sai em Cardo, que cobre politônico. A troca é automática — o
% pacote ucharclasses observa o bloco Unicode de cada caractere.
%
% Blocos cobertos:
%   Greek         U+0370–U+03FF  (grego moderno e básico)
%   GreekExtended U+1F00–U+1FFF  (politônico: espíritos e acentos)
%
% Usado por Novo-Caderno.ps1 quando a fonte escolhida não é a Cardo.
% Com Cardo como fonte principal, este arquivo é desnecessário.

\usepackage{fontspec}

% Scale=MatchLowercase iguala a altura-x da Cardo à da fonte principal,
% para o grego não parecer menor que o português ao redor.
\newfontfamily\gregofont{Cardo}[Scale=MatchLowercase]

\usepackage[Latin,Greek,GreekExtended]{ucharclasses}
\setTransitionsForGreek{\gregofont}{\normalfont}
